\chapter{Primes, Seen from Everywhere}

\begin{center}\itshape
The problem of distinguishing prime numbers from composite numbers and of resolving the latter into their prime factors is
known to be one of the most important and useful in arithmetic.\\
\upshape --- Carl Friedrich Gauss, \emph{Disquisitiones Arithmeticae} (1801), article 329
\end{center}

At the October 1903 meeting of the American Mathematical Society in New York, Frank Nelson Cole was announced to speak
``On the factorization of large numbers''. He walked to the blackboard without a word, and computed $2^{67}-1$ by hand:
$147573952589676412927$. Then he moved to a clean part of the board and multiplied, again by hand
\[
193707721\times761838257287.
\]
The two results agreed. Cole sat down. That, at least, is how E.~T.~Bell told the story in 1951: Cole had not said a single
word, and the audience applauded. The \emph{Bulletin} of the Society records the talk and its factorisation, but not the
silence or the applause, so those details rest on Bell alone. Bell also says that when he asked Cole, in 1911, how long the
factorisation had taken, the answer was ``three years of Sundays''.

There is a twist that makes the story better. Since Mersenne's list of 1644, $2^{67}-1$ had been believed to be prime. In
1876 Édouard Lucas had already proved that it is \emph{not}, with a test that never produces a factor. So for twenty-seven
years mathematicians knew that this number was composite and had no idea what its factors were. Primality and factoring
are different problems: one is easy, the other is believed to be hard, and modern cryptography lives in the gap between
them. This chapter looks at primes from every side the first five tomes provide: groups that detect them, binomial
coefficients that remember them, other rings where they split, and the probabilities that make testing them fast.

\begin{think}
How can one prove that a number is composite without finding a factor? Find a property that every prime has and check
that the number fails it. Which property of primes from Tome~1 would you use? And conversely: how could you convince a
sceptical friend, in a few minutes of computation, that a $100$-digit number is prime?
\end{think}

\section{Certificates and liars}

A primality test checks a property that every prime has. A composite number that has the property anyway is a
\emph{liar}, and the art is to find properties with few liars, or properties that come with a short proof.

\begin{exercise}[$\star$]\label{pb:fermatnum}
Let $F_n=2^{2^n}+1$ be the Fermat numbers.
\begin{enumerate}[label=(\alph*)]
\item Show that every prime factor of $F_n$, $n\ge2$, is congruent to $1$ modulo $2^{n+2}$.
\item Show that $641$ divides $F_5$, without dividing, using $641=5\cdot2^7+1=5^4+2^4$.
\item Show that two distinct Fermat numbers are coprime, and deduce that there are infinitely many primes.
\end{enumerate}
\end{exercise}

\begin{exercise}[$\star$]\label{pb:mersennediv}
Let $p$ be an odd prime and $M_p=2^p-1$.
\begin{enumerate}[label=(\alph*)]
\item Show that every prime factor $q$ of $M_p$ satisfies $q\equiv1\pmod{2p}$ and $q\equiv\pm1\pmod8$.
\item Show that if $p\equiv3\pmod4$ and $2p+1$ is prime, then $2p+1$ divides $M_p$. Deduce that $M_{11}$, $M_{23}$ and
$M_{83}$ are composite.
\end{enumerate}
\end{exercise}

\begin{exercise}[$\star\star$]\label{pb:pseudo}
An odd composite $n$ with $2^{n-1}\equiv1\pmod n$ is a \emph{pseudoprime} (to base $2$).
\begin{enumerate}[label=(\alph*)]
\item Show that $341=11\cdot31$ is a pseudoprime.
\item Show that if $n$ is a pseudoprime, so is $2^n-1$. Deduce that there are infinitely many pseudoprimes.
\item Show that $\frac{4^p-1}3$ is a pseudoprime for every prime $p\ge5$.
\end{enumerate}
\end{exercise}

\begin{exercise}[$\star\star$]\label{pb:carmichael}
A composite $n$ is a \emph{Carmichael number} if $a^{n-1}\equiv1\pmod n$ for every $a$ prime to $n$.
\begin{enumerate}[label=(\alph*)]
\item Show that a Carmichael number is odd, squarefree, and that $p-1\mid n-1$ for every prime $p\mid n$.
\item Show that a Carmichael number has at least three prime factors.
\item Show that if $6k+1$, $12k+1$ and $18k+1$ are all prime, their product is a Carmichael number. Find the smallest
example.
\end{enumerate}
\end{exercise}

\begin{exercise}[$\star\star$]\label{pb:euler}
Let $n$ be odd. Call $a$ an \emph{Euler liar} for $n$ if $\gcd(a,n)=1$ and $a^{(n-1)/2}\equiv\bigl(\frac an\bigr)\pmod n$, where
$\bigl(\frac an\bigr)$ is the Jacobi symbol.
\begin{enumerate}[label=(\alph*)]
\item Show that the Euler liars form a subgroup of $(\Z/n\Z)^*$.
\item Show that if $n$ is composite, this subgroup is proper, so that at most half of the units are Euler liars. In
particular, no odd composite number fools this test for all bases, unlike Problem~\ref{pb:carmichael}.
\end{enumerate}
\end{exercise}

\begin{exercise}[$\star\star$]\label{pb:aks}
Let $n\ge2$ and let $a$ be an integer prime to $n$. Show that $n$ is prime if and only if
\[
(X+a)^n\equiv X^n+a\pmod n
\]
in $\Z[X]$, that is, coefficient by coefficient.
\end{exercise}

\begin{exercise}[$\star\star$]\label{pb:pratt}
\begin{enumerate}[label=(\alph*)]
\item Show that $n\ge3$ is prime if and only if there is an $a$ with $a^{n-1}\equiv1\pmod n$ and $a^{(n-1)/q}\not\equiv1\pmod n$
for every prime $q\mid n-1$.
\item A \emph{certificate} for an odd prime $p$ is such an $a$, the factorisation of $p-1$, and a certificate for each odd
prime factor of $p-1$, and so on down. Let $f(p)$ be the number of primes in the whole tree, with $f(2)=1$ and
$f(p)=1+\sum_{q\mid p-1}f(q)$, the sum over the distinct primes $q$. Show that $f(p)\le2\log_2p-1$ for every odd prime $p$.
\item Deduce that a prime of $k$ binary digits has a certificate that can be checked with a number of operations on
$k$-digit numbers bounded by a polynomial in $k$.
\end{enumerate}
\end{exercise}

\begin{exercise}[$\star\star$]\label{pb:pocklington}
\begin{enumerate}[label=(\alph*)]
\item Let $n>1$, and let $F\mid n-1$. Suppose that $a^{n-1}\equiv1\pmod n$ and that $\gcd(a^{(n-1)/q}-1,n)=1$ for every prime
$q\mid F$. Show that every prime factor of $n$ is $\equiv1\pmod F$. Deduce that $n$ is prime if $F>\sqrt n$.
\item (Proth) Let $N=k2^m+1$ with $k$ odd and $k<2^m$. Show that if $a^{(N-1)/2}\equiv-1\pmod N$ for some $a$, then $N$ is
prime. Check that $3329=13\cdot2^8+1$ is prime with $a=3$.
\end{enumerate}
\end{exercise}

\begin{exercise}[Journey]\label{pb:rabin}
\emph{The Miller--Rabin test lies at most one time in four.} Let $n$ be an odd composite number, $n-1=2^sd$ with $d$ odd.
An $a$ prime to $n$ is a \emph{strong liar} if $a^d\equiv1\pmod n$ or $a^{2^jd}\equiv-1\pmod n$ for some $0\le j<s$. Write
$n=\prod_{i=1}^rp_i^{e_i}$ and $p_i-1=2^{\nu_i}m_i$ with $m_i$ odd, and let $\nu=\min_i\nu_i$.
\begin{enumerate}[label=(\alph*)]
\item In a cyclic group $C$ of even order $m$, with $-1$ its element of order $2$, show that $x^t=1$ has $\gcd(t,m)$
solutions and $x^t=-1$ has $\gcd(2t,m)-\gcd(t,m)$ solutions.
\item (Monier) Show that the number of strong liars is
\[
\prod_{i=1}^r\gcd(d,m_i)\cdot\Bigl(1+\frac{2^{r\nu}-1}{2^r-1}\Bigr).
\]
\item Deduce that at most $\frac14\varphi(n)$ of the units are strong liars, unless $n=9$. Check that $15$ attains $\frac14$.
\end{enumerate}
\end{exercise}

\begin{figure}[ht]
\centering
\fig{ch03-liars}
\caption{For each odd composite $n<1500$, the proportion of units that are Euler liars (yellow) and strong liars (blue). The
yellow dots stay below $\frac12$ (Problem~\ref{pb:euler}), the blue ones below $\frac14$ (Problem~\ref{pb:rabin}), except for
$n=9$, at $\frac13$. Most composite numbers have almost no liars at all: the bounds are worst cases.}
\label{fig:liars}
\end{figure}

\begin{exercise}[$\star\star\star$]\label{pb:lucaslehmer}
\emph{The Lucas--Lehmer test.} Let $s_0=4$ and $s_{k+1}=s_k^2-2$. Let $p$ be an odd prime and $M_p=2^p-1$. Show that $M_p$
is prime if and only if $M_p\mid s_{p-2}$.
\end{exercise}

\section{Binomials that remember primes}

Wilson's theorem says that $(p-1)!$ knows $p$ modulo $p$. The binomial coefficients, harmonic sums and Bernoulli numbers
know more: they see $p$ modulo $p^2$ and $p^3$. Here a fraction $\frac ab$ with $p\nmid b$ is reduced modulo $p^k$ by
inverting $b$, and ``$x\equiv y\pmod{p^k}$'' for such fractions means that $x-y$ is $p^k$ times a fraction of the same kind.

\begin{exercise}[$\star\star$]\label{pb:wolsten}
\emph{Wolstenholme's theorem.} Let $p\ge5$ be prime.
\begin{enumerate}[label=(\alph*)]
\item Show that $\sum_{k=1}^{p-1}\frac1{k^2}\equiv0\pmod p$.
\item Show that $1+\frac12+\dots+\frac1{p-1}\equiv0\pmod{p^2}$.
\item Show that $\binom{2p-1}{p-1}\equiv1\pmod{p^3}$, that is, $\binom{2p}p\equiv2\pmod{p^3}$.
\end{enumerate}
\end{exercise}

\begin{exercise}[$\star\star$]\label{pb:jacobsthal}
Let $p$ be prime and $a\ge b\ge0$. Show that $\binom{ap}{bp}\equiv\binom ab\pmod{p^2}$.
\end{exercise}

\begin{exercise}[$\star\star$]\label{pb:fermatquot}
For a prime $p$ and $p\nmid a$, the \emph{Fermat quotient} is $q_p(a)=\frac{a^{p-1}-1}p$, an integer.
\begin{enumerate}[label=(\alph*)]
\item Show that $q_p(ab)\equiv q_p(a)+q_p(b)\pmod p$, and $q_p(a+pt)\equiv q_p(a)-\frac ta\pmod p$.
\item Show that $2q_p(2)\equiv\sum_{k=1}^{p-1}\frac{(-1)^{k-1}}k\pmod p$ for odd $p$.
\item (Lerch) Let $W_p=\frac{(p-1)!+1}p$. Show that $\sum_{a=1}^{p-1}q_p(a)\equiv W_p\pmod p$.
\item Show that $2^{p-1}\equiv1\pmod{p^2}$ if and only if $2$ has the same order modulo $p$ and modulo $p^2$.
\end{enumerate}
\end{exercise}

\begin{exercise}[$\star\star$]\label{pb:morley}
Let $p$ be an odd prime. Show that
\[
\binom{p-1}{(p-1)/2}\equiv(-1)^{(p-1)/2}\,4^{p-1}\pmod{p^2}.
\]
\end{exercise}

\begin{exercise}[$\star\star\star$]\label{pb:glaisher}
Let $p$ be an odd prime and $B_{p-1}$ the Bernoulli number. Show that $pB_{p-1}$ has no $p$ in its denominator and that
\[
(p-1)!\equiv pB_{p-1}-p\pmod{p^2}.
\]
\end{exercise}

\section{Primes in other places}

A prime of $\Z$ may split in a larger ring, and a sequence built from primes may avoid some of them forever. Counting needs
the prime number theorem and its relatives from Tome~5.

\begin{exercise}[$\star\star$]\label{pb:gaussprimes}
Let $G(x)$ be the number of Gaussian primes $\pi\in\Z[i]$ with $N(\pi)\le x$, all associates counted.
\begin{enumerate}[label=(\alph*)]
\item Show that $G(x)=4\bigl(1+2\pi(x;4,1)+\pi(\sqrt x;4,3)\bigr)$, where $\pi(y;4,a)$ counts the primes $p\le y$ with
$p\equiv a\pmod4$.
\item Show that $G(x)\sim\frac{4x}{\log x}$. Since about $\pi x$ Gaussian integers have norm at most $x$, what is the
``probability'' that a Gaussian integer of norm near $N$ is prime?
\item Show that $\sum\frac1{N(\pi)}=\log\log x+O(1)$, the sum over the Gaussian primes of norm at most $x$, one in each
class of associates.
\end{enumerate}
\end{exercise}

\begin{figure}[ht]
\centering
\fig{ch03-gaussprimes}
\caption{The Gaussian primes of norm at most $22^2$. The red ones lie on the axes (the primes $p\equiv3\pmod4$ of $\Z$, and
their associates) or have norm $2$; the blue ones come in groups of eight, $\pm a\pm bi$ and $\pm b\pm ai$, one group for each
prime $p=a^2+b^2\equiv1\pmod4$. The picture has the symmetry of the square, and no other.}
\label{fig:gaussprimes}
\end{figure}

\begin{exercise}[$\star\star$]\label{pb:mullin}
\begin{enumerate}[label=(\alph*)]
\item Let $a_1=2$ and let $a_{n+1}$ be the smallest prime factor of $a_1a_2\cdots a_n+1$. Show that the $a_n$ are distinct, and
compute $a_1,\dots,a_7$.
\item Let $b_1=2$ and let $b_{n+1}$ be the \emph{largest} prime factor of $b_1\cdots b_n+1$. Show that $5$ never appears.
\end{enumerate}
\end{exercise}

\begin{exercise}[$\star\star\star$]\label{pb:nsq1}
Let $P(x)$ be the largest prime factor of $\prod_{1\le n\le x}(n^2+1)$. Show that $P(x)>x$ for all large $x$.
\end{exercise}

\begin{exercise}[$\star\star$]\label{pb:nair}
\begin{enumerate}[label=(\alph*)]
\item Let $d_n=\operatorname{lcm}(1,2,\dots,n)$. Show that $m\binom nm$ divides $d_n$ for $1\le m\le n$.
\item Deduce that $d_n\ge2^n$ for $n\ge7$, and that $\pi(n)\ge\frac{n\log2}{\log n}$ for $n\ge7$.
\end{enumerate}
\end{exercise}

\section{Primes and chance}

The tests of the first section are randomised: they pick bases at random and bet. How long does one wait for a prime,
and how small is the first base that a prime cannot fool?

\begin{exercise}[$\star$]\label{pb:artin}
Let $q$ be a prime such that $p=2q+1$ is also prime, with $p>7$. Show that $2$ is a primitive root modulo $p$ if
$q\equiv1\pmod4$, and that $-2$ is one if $q\equiv3\pmod4$.
\end{exercise}

\begin{exercise}[$\star$]\label{pb:findprime}
To make an RSA key, one draws random integers in $[N,2N]$ until one is prime.
\begin{enumerate}[label=(\alph*)]
\item Show that the expected number of draws is asymptotic to $\log N$.
\item If one only draws integers prime to $30$, show that it becomes $\sim\frac4{15}\log N$. With only integers prime to all
primes $\le y$ ($y$ fixed, then large), what is the saving?
\item Estimate both numbers for $N=2^{1023}$.
\end{enumerate}
\end{exercise}

\begin{exercise}[$\star\star$]\label{pb:leastqnr}
Let $p$ be an odd prime and $n_p$ the least positive quadratic non-residue modulo $p$.
\begin{enumerate}[label=(\alph*)]
\item Show that $n_p$ is a prime and that $n_p<\sqrt p+1$.
\item Pretend that the symbols $\bigl(\frac qp\bigr)$, $q=2,3,5,7,\dots$ prime, are independent fair signs. Show that the expected
value of $n_p$ in this model is $\sum_{k\ge1}p_k2^{-k}$, where $p_k$ is the $k$-th prime, and compute it to three decimals.
\end{enumerate}
\end{exercise}

\section{Hints}

\begin{hint}{pb:fermatnum}
(a) The order of $2$ modulo $p$ is $2^{n+1}$, and $2$ is a square modulo $p$ when $p\equiv1\pmod8$. (c) $F_0F_1\cdots F_{n-1}=F_n-2$.
\end{hint}

\begin{hint}{pb:mersennediv}
(a) The order of $2$ modulo $q$ is $p$, and $2=2^{p+1}$ is a square modulo $q$. (b) $2p+1\equiv7\pmod8$.
\end{hint}

\begin{hint}{pb:pseudo}
(b) $2^n-1$ is composite when $n$ is, and $2^n\equiv1$ modulo $2^n-1$; how does $n$ divide $2^n-2$? (c) $3N=2^{2p}-1$.
\end{hint}

\begin{hint}{pb:carmichael}
(a) Try $a=-1$; then use a primitive root modulo $p^2$ or modulo $p$, glued by the Chinese remainder theorem. (b) If $n=pq$
with $p<q$, compute $n-1$ modulo $q-1$.
\end{hint}

\begin{hint}{pb:euler}
(b) If $p^2\mid n$, use an element of order $p(p-1)$ modulo $p^2$. If $n=pm$ with $p\nmid m$, take $a$ a non-residue modulo $p$
and $a\equiv1\pmod m$.
\end{hint}

\begin{hint}{pb:aks}
If $q^k$ is the exact power of a prime $q$ dividing $n$, look at the coefficient of $X^q$.
\end{hint}

\begin{hint}{pb:pratt}
(a) An element of order $n-1$ forces $\varphi(n)=n-1$. (b) Induction, using $\prod q_i\le\frac{p-1}2$ over the odd primes $q_i\mid p-1$.
\end{hint}

\begin{hint}{pb:pocklington}
(a) For a prime $p\mid n$, the order of $a$ modulo $p$ divides $n-1$ but not $(n-1)/q$. (b) $F=2^m$.
\end{hint}

\begin{hint}{pb:rabin}
(b) Use the Chinese remainder theorem and the cyclic groups $(\Z/p_i^{e_i}\Z)^*$; note that $p_i\nmid d$. (c) Compare with
$\varphi(n)=\prod p_i^{e_i-1}2^{\nu_i}m_i$. For $r=2$, the hardest case is $n=p_1p_2$ with $\nu_1=\nu_2$.
\end{hint}

\begin{hint}{pb:lucaslehmer}
$s_k=\omega^{2^k}+\bar\omega^{2^k}$ with $\omega=2+\sqrt3$, $\omega\bar\omega=1$. One direction: the order of $\omega$ in
$(\Z/q\Z)[\sqrt3]$. The other: $\omega=\frac{(1+\sqrt3)^2}2$ and Frobenius in $\mathbb F_{M_p^2}$.
\end{hint}

\begin{hint}{pb:wolsten}
(a) $k\mapsto\frac1k$ permutes the non-zero residues. (b) Pair $k$ with $p-k$. (c) $\binom{2p-1}{p-1}=\prod_{k=1}^{p-1}\bigl(1+\frac pk\bigr)$.
\end{hint}

\begin{hint}{pb:jacobsthal}
Vandermonde with $ap=(a-1)p+p$, and induction on $a$.
\end{hint}

\begin{hint}{pb:fermatquot}
(b) Expand $2^p=(1+1)^p$; $\frac1p\binom pk\equiv\frac{(-1)^{k-1}}k\pmod p$. (c) Apply (a) to $(p-1)!$.
\end{hint}

\begin{hint}{pb:morley}
$\binom{p-1}k=\prod_{j\le k}\bigl(\frac pj-1\bigr)$, and relate $H_{(p-1)/2}$ to $q_p(2)$ through Problem~\ref{pb:fermatquot}(b).
\end{hint}

\begin{hint}{pb:glaisher}
Faulhaber's formula for $\sum_{k<p}k^{p-1}$ modulo $p^2$, von Staudt--Clausen, and Problem~\ref{pb:fermatquot}(c).
\end{hint}

\begin{hint}{pb:gaussprimes}
(a) The Gaussian primes up to units are $1+i$, the $\pi$ and $\bar\pi$ above each $p\equiv1\pmod4$, and the $p\equiv3\pmod4$. (b)
and (c): the prime number theorem and Dirichlet's theorem in progressions.
\end{hint}

\begin{hint}{pb:mullin}
(b) What would $b_1\cdots b_n+1$ be, if its largest prime factor were $5$? Look at it modulo $4$.
\end{hint}

\begin{hint}{pb:nsq1}
Compare $\log\prod(n^2+1)\ge2\log\lfloor x\rfloor!$ with what primes $p\le x$, $p\equiv1\pmod4$, can contribute: at most two
solutions of $n^2\equiv-1\pmod{p^k}$.
\end{hint}

\begin{hint}{pb:nair}
(a) $\int_0^1x^{m-1}(1-x)^{n-m}\dd x$ in two ways. (b) $n=2k+1$ and $m=k+1$, then $m=k$ with $n=2k$.
\end{hint}

\begin{hint}{pb:artin}
The order of $2$ divides $2q$; Euler's criterion and the value of $\bigl(\frac2p\bigr)$ decide whether it is $q$.
\end{hint}

\begin{hint}{pb:findprime}
A geometric number of draws. (b) Mertens's product $\prod_{p\le y}(1-\frac1p)\sim\frac{e^{-\gamma}}{\log y}$.
\end{hint}

\begin{hint}{pb:leastqnr}
(a) Look at $m=\lfloor p/n_p\rfloor+1$ and $mn_p-p$.
\end{hint}

\section{Solutions}

\begin{solution}{pb:fermatnum}
\idea A prime factor $p$ of $F_n$ makes $2$ have a large order, a power of $2$, which must divide $p-1$. That already gives
$2^{n+1}\mid p-1$; the extra factor $2$ comes from knowing that $2$ is a square. For $641$, the two expressions of $641$ say
that $5\cdot2^7\equiv-1$ and $2^4\equiv-5^4$, and multiplying the right powers gives $2^{32}\equiv-1$.

\writeup (a) Let $p\mid F_n$. Then $2^{2^n}\equiv-1$, so $2^{2^{n+1}}\equiv1$ and the order of $2$ modulo $p$ divides $2^{n+1}$ but
not $2^n$: it is $2^{n+1}$, and it divides $p-1$ (Tome~1, Chapter~5). For $n\ge2$ this gives $p\equiv1\pmod8$, so
$\bigl(\frac2p\bigr)=1$ (Tome~1, Chapter~6) and, by Euler's criterion, $2^{(p-1)/2}\equiv1$. So $2^{n+1}\mid\frac{p-1}2$, and
$p\equiv1\pmod{2^{n+2}}$.

(b) Modulo $641$: $5\cdot2^7\equiv-1$, so $5^4\cdot2^{28}\equiv1$; and $2^4\equiv-5^4$. Hence
$2^{32}=2^4\cdot2^{28}\equiv-5^4\cdot2^{28}\equiv-1$, and $641\mid F_5$. In accordance with (a), $641=10\cdot2^6+1\equiv1\pmod{2^7}$: Euler
needed only a few trials.

(c) $F_0F_1\cdots F_{n-1}=F_n-2$, by induction: $(F_n-2)F_n=(2^{2^n}-1)(2^{2^n}+1)=2^{2^{n+1}}-1=F_{n+1}-2$. So a common divisor
of $F_m$ and $F_n$, $m<n$, divides $2$; the $F_n$ are odd, so they are coprime. Each $F_n$ has a prime factor, and these are
all different: there are infinitely many primes.
\tools Tome~1, Chapters~5 and~6 (orders, the value of $(\frac2p)$, Euler's criterion).
\leads Fermat conjectured in 1640 that all $F_n$ are prime; Euler found $641$ in 1732. Only $F_0,\dots,F_4$ are known to be prime,
and $F_5,\dots,F_{32}$ are known to be composite. Landry found $F_6=274177\cdot67280421310721$ in 1880, and indeed
$274177\equiv1\pmod{2^8}$. Gauss proved that the regular $n$-gon is constructible when $n$ is a power of $2$ times distinct Fermat
primes (Tome~2, Chapter~10).
\end{solution}

\begin{solution}{pb:mersennediv}
\idea Same game as for Fermat numbers. The order of $2$ modulo a factor $q$ is the prime $p$, so $p\mid q-1$. And
$2^{p+1}=2\cdot2^p\equiv2$ exhibits $2$ as a square, which pins $q$ modulo $8$. For (b), Euler's criterion turns
``$2$ is a square modulo $2p+1$'' into ``$2^p\equiv1\pmod{2p+1}$''.

\writeup (a) $2^p\equiv1\pmod q$ and $2\not\equiv1$, so the order of $2$ is $p$, which divides $q-1$; as $q$ is odd, $2p\mid q-1$.
Moreover $p+1$ is even and $2\equiv2^{p+1}=\bigl(2^{(p+1)/2}\bigr)^2$, so $\bigl(\frac2q\bigr)=1$ and $q\equiv\pm1\pmod8$
(Tome~1, Chapter~6).

(b) If $p\equiv3\pmod4$, $q=2p+1\equiv7\pmod8$, so $\bigl(\frac2q\bigr)=1$ and Euler's criterion gives $2^{(q-1)/2}=2^p\equiv1\pmod q$. So
$q\mid M_p$, and $q<M_p$ when $p>3$, so $M_p$ is composite: $23\mid M_{11}=2047$, $47\mid M_{23}$, $167\mid M_{83}$.
\tools Tome~1, Chapters~5 and~6.
\leads Euler stated (b) in 1750. By (a), the trial divisions for $M_p$ need only the primes $2kp+1\equiv\pm1\pmod8$, which is
how the GIMPS project eliminates most exponents before running the Lucas--Lehmer test of Problem~\ref{pb:lucaslehmer}. Cole's
factor $193707721=2\cdot1445580\cdot67+1$ is of the form of (a).
\end{solution}

\begin{solution}{pb:pseudo}
\idea For (b), the number $m=2^n-1$ satisfies $2^n\equiv1\pmod m$, so $2^{m-1}\equiv1$ as soon as $n\mid m-1$. And
$m-1=2(2^{n-1}-1)$, which $n$ divides exactly because $n$ is a pseudoprime. For (c), the same mechanism with $2^{2p}$.

\writeup (a) $2^{10}=1024=93\cdot11+1$ and $2^5=32=31+1$, so $2^{10}\equiv1\pmod{11}$ and $2^5\equiv1\pmod{31}$. Since $10\mid340$ and
$5\mid340$, $2^{340}\equiv1$ modulo $11$ and $31$, hence modulo $341$.

(b) Let $n$ be a pseudoprime and $m=2^n-1$, which is odd. If $n=ab$ with $1<a<n$, then $2^a-1$ divides $2^n-1$ and lies strictly
between $1$ and $m$, so $m$ is composite. Now $m-1=2(2^{n-1}-1)$ and $n\mid2^{n-1}-1$, so $m-1=nk$ for an integer $k$. As
$2^n\equiv1\pmod m$, $2^{m-1}=(2^n)^k\equiv1\pmod m$. So $m$ is a pseudoprime, larger than $n$. Starting from $341$, this gives
infinitely many.

(c) Let $N=\frac{4^p-1}3=(2^p-1)\cdot\frac{2^p+1}3$. Both factors are integers
larger than $1$ ($3\mid2^p+1$ as $p$ is odd, and $\frac{2^p+1}3>1$ for $p\ge3$). So $N$ is composite, and odd. Next,
$N-1=\frac{4^p-4}3=4\cdot\frac{4^{p-1}-1}3$. By Fermat, $p\mid4^{p-1}-1$, and as $p\ne3$, $p$ divides $\frac{4^{p-1}-1}3$. So $2p\mid N-1$.
Finally $3N=2^{2p}-1$, so $2^{2p}\equiv1\pmod N$ and $2^{N-1}\equiv1\pmod N$.
\tools Tome~1, Chapters~1 and~3 (divisibility, Fermat's little theorem).
\leads (b) is due to Malo (1903) and (c) to Cipolla (1904). Pseudoprimes are rare: Pomerance proved that their number up to $x$
is at most $x\,e^{-\frac12\log x\log\log\log x/\log\log x}$ for large $x$, so that a random odd number that passes the base $2$
test is almost surely prime, which is what software relied on for decades.
\end{solution}

\begin{solution}{pb:carmichael}
\idea To show that $n$ fails, find one base $a$ that breaks $a^{n-1}\equiv1$. The Chinese remainder theorem lets us choose
$a$ freely modulo each prime power factor, so we can take it of maximal order there. Then $a^{n-1}\equiv1$ forces that order
to divide $n-1$. With two prime factors, the divisibilities $p-1\mid n-1$ and $q-1\mid n-1$ cannot both hold.

\writeup (a) $n$ is odd: otherwise $(-1)^{n-1}=-1\not\equiv1\pmod n$, as $n>2$. Let $p\mid n$ and write $n=p^vm$ with $p\nmid m$.
If $v\ge2$, let $g$ have order $p(p-1)$ modulo $p^2$ ($(\Z/p^2\Z)^*$ is cyclic, Tome~1, Chapter~5), and by the Chinese
remainder theorem choose $a\equiv g\pmod{p^2}$, $a\equiv1\pmod m$, and $a$ prime to $p$; then $a$ is prime to $n$. From
$a^{n-1}\equiv1\pmod{p^2}$, $p(p-1)\mid n-1$, so $p\mid n-1$, absurd since $p\mid n$. So $n$ is squarefree. Now take $a\equiv g$
modulo $p$ with $g$ a primitive root, and $a\equiv1\pmod m$: $a^{n-1}\equiv1\pmod p$ gives $p-1\mid n-1$.

(b) $n$ is squarefree and composite, so it has at least two prime factors. If $n=pq$ with $p<q$, then
$n-1=p(q-1)+(p-1)\equiv p-1\pmod{q-1}$, and $q-1\mid n-1$ would give $q-1\mid p-1$, impossible as $0<p-1<q-1$.

(c) $n=(6k+1)(12k+1)(18k+1)=1296k^3+396k^2+36k+1$, so $n-1=36k(36k^2+11k+1)$ is divisible by $6k$, $12k$ and $18k$. By
Korselt's criterion (Tome~1, Chapter~3), $n$ is a Carmichael number. For $k=1$: $7\cdot13\cdot19=1729$.
\tools Tome~1, Chapters~3 and~5 (Chinese remainder theorem, Korselt's criterion, primitive roots modulo $p^2$).
\leads $1729$ is Hardy's taxi-cab number, Ramanujan's ``smallest number expressible as a sum of two cubes in two different
ways''. Chernick (1939) gave (c); whether it gives infinitely many examples is open, but Alford, Granville and Pomerance proved
in 1994 that there are infinitely many Carmichael numbers.
\end{solution}

\begin{solution}{pb:euler}
\idea Both sides of the congruence are multiplicative in $a$, so the liars form a subgroup; a proper subgroup has index at
least $2$. To show it is proper, we need one unit that is not a liar, and as for Carmichael numbers we build it with the
Chinese remainder theorem. The Jacobi symbol sees each prime factor separately, so one can make it $-1$ while $a^{(n-1)/2}$
stays $1$ modulo a part of $n$.

\writeup (a) If $a,b$ are liars, $(ab)^{(n-1)/2}\equiv\bigl(\frac an\bigr)\bigl(\frac bn\bigr)=\bigl(\frac{ab}n\bigr)$ by
multiplicativity of the Jacobi symbol in $a$ (Tome~1, Chapter~6). The set contains $1$ and is finite, so it is a subgroup.

(b) Let $n$ be odd and composite. If $p^2\mid n$ for some $p$, choose $a$ of order $p(p-1)$ modulo $p^2$ and $\equiv1$ modulo the rest
of $n$, as in Problem~\ref{pb:carmichael}. A liar satisfies $a^{n-1}\equiv\bigl(\frac an\bigr)^2=1\pmod n$, so $p(p-1)\mid n-1$,
impossible. Otherwise $n=pm$ with $m>1$ odd and $p\nmid m$. Let $b$ be a non-residue modulo $p$, and choose $a\equiv b\pmod p$,
$a\equiv1\pmod m$. Then $\bigl(\frac an\bigr)=\bigl(\frac ap\bigr)\bigl(\frac am\bigr)=\bigl(\frac bp\bigr)\bigl(\frac1m\bigr)=-1$. If $a$
were a liar, $a^{(n-1)/2}\equiv-1\pmod n$, hence modulo $m$; but $a\equiv1\pmod m$, so $1\equiv-1\pmod m$ and $m\mid2$: absurd.

A proper subgroup of a finite group has index at least $2$ (Lagrange, Tome~2, Chapter~1), so at most half of the units are
liars.
\tools Tome~1, Chapters~3, 5 and~6 (Chinese remainder theorem, primitive roots, Jacobi symbol); Tome~2, Chapter~1 (Lagrange).
\leads This is the Solovay--Strassen test (1977), the first randomised primality test: with $k$ random bases, a composite
passes with probability at most $2^{-k}$. Its proof of correctness is what made ``probably prime'' a mathematical statement.
The Miller--Rabin test of Problem~\ref{pb:rabin} is faster and lies less.
\end{solution}

\begin{solution}{pb:aks}
\idea Expand with the binomial theorem: the identity holds if and only if $n$ divides all the binomial coefficients $\binom nk$,
$0<k<n$, and $a^n-a$. For a prime this is Tome~1. For a composite, one coefficient escapes, and the natural candidate is
$\binom nq$ for a prime factor $q$: its $q$-adic valuation is easy to compute.

\writeup Expand:
\[
(X+a)^n-(X^n+a)=\sum_{0<k<n}\binom nka^{n-k}X^k+(a^n-a).
\]
If $n$ is prime, $n\mid\binom nk$ for $0<k<n$ and $n\mid a^n-a$
(Tome~1, Chapter~3), so the congruence holds.

Let $n$ be composite, $q$ a prime factor, $q^k$ the exact power of $q$ in $n$; then $q<n$. The coefficient of $X^q$ is
$\binom nqa^{n-q}$. In $\binom nq=\frac{n(n-1)\cdots(n-q+1)}{q!}$, the numerator contains $q$ consecutive integers, exactly one of which,
$n$, is divisible by $q$, so its $q$-adic valuation is $k$; and $v_q(q!)=1$. So $v_q\bigl(\binom nq\bigr)=k-1$. Since $q\nmid a$, $q^k$
does not divide the coefficient of $X^q$, nor does $n$. The congruence fails.
\tools Tome~1, Chapters~1 and~3 (valuations, Fermat's little theorem, binomial coefficients modulo $p$).
\leads The criterion is useless as it stands: it has $n$ coefficients. Agrawal, Kayal and Saxena (2002) showed that it is
enough to check it modulo $X^r-1$ for one small $r$ and for few values of $a$, with $r$ and the number of $a$'s polynomial in
$\log n$. This gave the first deterministic polynomial-time primality test, ``PRIMES is in P''.
\end{solution}

\begin{solution}{pb:pratt}
\idea (a) An element whose order is exactly $n-1$ forces the group $(\Z/n\Z)^*$ to have $n-1$ elements, which only happens
for primes; and checking that the order is $n-1$ needs only the prime factors of $n-1$. (b) The tree branches at the prime
factors of $p-1$, whose product is less than $p$: induction on $p$ with a logarithmic bound, because logarithms add.

\writeup (a) If $n$ is prime, a primitive root works (Tome~1, Chapter~5). Conversely, the order $e$ of $a$ in $(\Z/n\Z)^*$ ($a$ is a
unit since $a\cdot a^{n-2}\equiv1$) divides $n-1$. If $e<n-1$, some prime $q$ divides $\frac{n-1}e$, and then $e\mid\frac{n-1}q$ and
$a^{(n-1)/q}\equiv1$, excluded. So $e=n-1$, hence $n-1\le\varphi(n)$, and $n$ is prime.

(b) Induction on $p$. As $p$ is odd, $2\mid p-1$. Let $q_1,\dots,q_k$ be the distinct odd primes dividing $p-1$; their product is at
most $\frac{p-1}2$. If $k=0$, $f(p)=1+f(2)=2\le2\log_2p-1$, because $p\ge3$. If $k\ge1$, by induction
\[
f(p)=1+f(2)+\sum_{i=1}^kf(q_i)\le2+\sum_i(2\log_2q_i-1)\le2+2\log_2\frac{p-1}2-k=2\log_2(p-1)-k,
\]
which is at most $2\log_2p-1$.

(c) A prime $p$ with $k$ digits has at most $2k$ primes in its tree, each with at most $k$ digits. At each one, say $\ell$, the
checker multiplies the claimed factorisation of $\ell-1$ (at most $k$ factors with multiplicity) and computes $a^{\ell-1}$ and the
$a^{(\ell-1)/q}$ modulo $\ell$, at most $k+1$ modular powers, each by $O(k)$ squarings and multiplications (Tome~1, Chapter~3). In
total $O(k^3)$ operations on numbers of $k$ digits.
\tools Tome~1, Chapters~3 and~5 (fast exponentiation, orders, primitive roots).
\leads Pratt (1975): \emph{PRIMES is in NP}, the first complexity result on primality. Finding the certificate requires factoring
$p-1$, which is hard in general; but primes of special shape (Proth primes, Problem~\ref{pb:pocklington}) carry their own
factorisation, and elliptic curve methods replace $p-1$ by the order of a random group (Goldwasser--Kilian, Atkin--Morain).
\end{solution}

\begin{solution}{pb:pocklington}
\idea Lucas's criterion asks to factor all of $n-1$. Pocklington noticed that a partial factorisation is enough: every prime
factor of $n$ is then $\equiv1$ modulo the factored part $F$, hence larger than $F$, and a composite $n$ has a prime factor at
most $\sqrt n$. For Proth numbers $N-1=k2^m$, the power of $2$ is the factored part.

\writeup (a) Let $p$ be a prime factor of $n$ and $e$ the order of $a$ modulo $p$ ($p\nmid a$ since $a^{n-1}\equiv1$). Then
$e\mid n-1$. Let $q^f$ be the exact power of a prime $q\mid F$ in $F$, and $q^g$ its power in $n-1$, $g\ge f$. If $q^g\nmid e$, then
$e\mid\frac{n-1}q$ and $p\mid a^{(n-1)/q}-1$, contradicting $\gcd(a^{(n-1)/q}-1,n)=1$. So $q^g\mid e\mid p-1$. Hence $F\mid p-1$ and
$p>F$. If $F>\sqrt n$, every prime factor of $n$ exceeds $\sqrt n$, so $n$ is prime.

(b) Take $F=2^m$, with only $q=2$. From $a^{(N-1)/2}\equiv-1$: $a^{N-1}\equiv1$, and $\gcd(a^{(N-1)/2}-1,N)=\gcd(-2,N)=1$ as $N$ is odd. And
$N=k2^m+1\le(2^m-1)2^m+1<2^{2m}$, so $F>\sqrt N$. By (a), $N$ is prime. For $3329$: $m=8$, $k=13<256$, and repeated squaring
gives $3^{1664}\equiv3328\equiv-1\pmod{3329}$, so $3329$ is prime.
\tools Tome~1, Chapters~3 and~5 (orders and fast exponentiation).
\leads Pocklington (1914) and Proth (1878). The largest known primes that are not Mersenne primes are Proth primes, found by the
PrimeGrid project with exactly this test. Note the asymmetry: (b) proves primality in one exponentiation, but a failure says
nothing unless $\bigl(\frac aN\bigr)=-1$.
\end{solution}

\begin{solution}{pb:rabin}
\idea A strong liar is described componentwise by the Chinese remainder theorem, and in each cyclic group $(\Z/p_i^{e_i}\Z)^*$
the equations $x^t=\pm1$ have explicit numbers of solutions. So the liars can be counted exactly, and the count is then
compared with $\varphi(n)$ factor by factor. The prime powers lose a factor $p^{e-1}$; several primes lose a factor
$2^{r-1}$ from the powers of $2$; and the only delicate case, two primes with the same $2$-part, is saved by a
divisibility argument.

\writeup (a) Let $g$ generate $C$, of order $m$. $x=g^u$ satisfies $x^t=1$ iff $m\mid tu$, iff $\frac m{\gcd(t,m)}\mid u$: there are
$\gcd(t,m)$ such $u$ modulo $m$. If $x^t=-1$, then $x^{2t}=1$ and $x^t\ne1$. Conversely if $x^{2t}=1\ne x^t$, then $x^t$ has order $2$,
so it is $-1$, the unique element of order $2$ in a cyclic group of even order. So the count is $\gcd(2t,m)-\gcd(t,m)$.

(b) Let $\varphi_i=p_i^{e_i-1}2^{\nu_i}m_i$, the order of the cyclic group $C_i=(\Z/p_i^{e_i}\Z)^*$ (Tome~1, Chapter~5; $p_i$ is odd).
By the Chinese remainder theorem, $a$ is a liar iff either $a^d=1$ in every $C_i$, or for some $0\le j<s$, $a^{2^jd}=-1$ in every
$C_i$. These cases are disjoint: if $a^{2^jd}\equiv-1$, then $a^{2^{j'}d}\equiv1$ for $j'>j$, and $\pm1$ are distinct modulo $n$. Since
$p_i\mid n$, $p_i\nmid n-1$, so $p_i\nmid d$, and $\gcd(2^jd,\varphi_i)=g_i\,2^{\min(j,\nu_i)}$ with $g_i=\gcd(d,m_i)$ ($d$ is odd). By (a):
\begin{itemize}
\item $a^d=1$ in $C_i$ for $g_i$ elements;
\item $a^{2^jd}=-1$ in $C_i$ for $g_i\bigl(2^{\min(j+1,\nu_i)}-2^{\min(j,\nu_i)}\bigr)$ elements, which is $g_i2^j$ if $j<\nu_i$ and $0$ if
$j\ge\nu_i$.
\end{itemize}
Multiplying over $i$, the number of $a$ with $a^{2^jd}\equiv-1\pmod n$ is $\bigl(\prod_ig_i\bigr)2^{jr}$ if $j<\nu$ and $0$ otherwise.
Each $p_i\equiv1\pmod{2^\nu}$, so $n\equiv1\pmod{2^\nu}$ and $\nu\le s$. Summing over $0\le j<\nu$ and adding the first case, the
number of liars is $\prod_ig_i\cdot\bigl(1+\sum_{j<\nu}2^{jr}\bigr)$, which is Monier's formula.

(c) Write $T_r(\nu)=\frac{1+\frac{2^{r\nu}-1}{2^r-1}}{2^{r\nu}}=2^{-r\nu}\Bigl(1-\frac1{2^r-1}\Bigr)+\frac1{2^r-1}$. Then the proportion of liars is
\[
R=\frac{L}{\varphi(n)}=\prod_{i=1}^r\frac{g_i}{m_ip_i^{e_i-1}}\cdot T_r(\nu)\cdot2^{r\nu-(\nu_1+\dots+\nu_r)}.
\]
All three factors are at most $1$, and $T_r$ decreases in $\nu\ge1$, so $T_r(\nu)\le T_r(1)=2^{-r}+2^{-r}=2^{1-r}$ (for $r=1$,
$T_1(\nu)=1$).

\emph{$r=1$}, $n=p^e$ with $e\ge2$: $R\le p^{1-e}$, which is $\frac13$ for $n=9$ and at most $\frac15$ otherwise ($p^{e-1}\ge5$).

\emph{$r\ge3$}: $R\le2^{1-r}\le\frac14$.

\emph{$r=2$}: $R\le\frac12\cdot\frac{g_1g_2}{m_1m_2p_1^{e_1-1}p_2^{e_2-1}}\cdot2^{2\nu-\nu_1-\nu_2}$. If $\nu_1\ne\nu_2$, the last factor is at most
$\frac12$. If some $e_i\ge2$, the middle factor is at most $\frac13$. There remains $n=p_1p_2$ with $\nu_1=\nu_2=\nu$. If
$g_1=m_1$ and $g_2=m_2$, then $m_i\mid d$, so $p_i-1=2^\nu m_i$ divides $2^sd=n-1$. But $n-1=p_1(p_2-1)+(p_1-1)$, so $p_2-1\mid p_1-1$,
and symmetrically $p_1-1\mid p_2-1$, so $p_1=p_2$: absurd. So some $g_i$ is a proper divisor of the odd $m_i$, $\frac{g_i}{m_i}\le\frac13$,
and $R\le\frac16$.

In all cases $R\le\frac14$, except $n=9$. For $n=15$: $14=2\cdot7$, $s=1$, $d=7$; $r=2$, $p_1-1=2$, $p_2-1=4$, $\nu=1$, $g_1=g_2=1$. Monier
gives $1\cdot(1+1)=2$ liars ($\pm1$) out of $\varphi(15)=8$: exactly $\frac14$.
\tools Tome~1, Chapters~3 and~5 (Chinese remainder theorem, cyclic groups $(\Z/p^e\Z)^*$); Tome~2, Chapter~1 (cyclic groups).
\leads Miller (1976) proved that under the generalised Riemann hypothesis the bases $a<2\log^2n$ suffice, a deterministic test;
Rabin (1980) turned it into the randomised test with error $4^{-k}$ for $k$ bases. Figure~\ref{fig:liars} shows that the
bound is far from typical: for most $n$ the proportion of liars is tiny, which is why, in practice, a handful of bases is
used for numbers of thousands of digits.
\end{solution}

\begin{solution}{pb:lucaslehmer}
\idea The recursion $s\mapsto s^2-2$ is the doubling formula for $x+x^{-1}$: if $s=x+x^{-1}$ then $s^2-2=x^2+x^{-2}$. With
$x=\omega=2+\sqrt3$, $s_k=\omega^{2^k}+\omega^{-2^k}$, and $M_p\mid s_{p-2}$ says that $\omega^{2^{p-1}}=-1$, so $\omega$ has order $2^p$,
modulo $M_p$. If $M_p$ had a small prime factor $q$, the group where $\omega$ lives modulo $q$ would be too small for that order.
Conversely, if $M_p$ is prime, the Frobenius of $\mathbb F_{M_p^2}$ computes $\omega^{(M_p+1)/2}$, once $\omega$ is written as a square.

\writeup Let $M=M_p$, $\omega=2+\sqrt3$, $\bar\omega=2-\sqrt3$; $\omega\bar\omega=1$. By induction $s_k=\omega^{2^k}+\bar\omega^{2^k}$:
$s_0=4$ and $(\omega^{2^k}+\bar\omega^{2^k})^2-2=\omega^{2^{k+1}}+\bar\omega^{2^{k+1}}$.

\emph{If $M\mid s_{p-2}$, then $M$ is prime.} Let $q$ be a prime factor of $M$ ($q\ne2$, and $q\ne3$ since $2^p\equiv2\pmod3$ makes
$M\equiv1\pmod3$). Work in the finite ring $R=(\Z/q\Z)[X]/(X^2-3)$, with $q^2$ elements, where $\sqrt3$ is the class of $X$. In $R$,
$\omega^{2^{p-2}}+\bar\omega^{2^{p-2}}=0$; multiply by $\omega^{2^{p-2}}$: $\omega^{2^{p-1}}=-1$, so $\omega^{2^p}=1$ and $\omega$ is a unit of order
$2^p$ (as $-1\ne1$ in $R$). The units of $R$ form a group with at most $q^2-1$ elements, so $2^p\le q^2-1$. If $M$ were composite, its
smallest prime factor would satisfy $q^2\le M=2^p-1$, and $2^p\le2^p-2$: absurd.

\emph{If $M$ is prime, then $M\mid s_{p-2}$.} As $p\ge3$, $M\equiv7\pmod8$, so $\bigl(\frac2M\bigr)=1$. And $M\equiv1\pmod3$, $M\equiv3\pmod4$,
so by reciprocity $\bigl(\frac3M\bigr)=-\bigl(\frac M3\bigr)=-1$ (Tome~1, Chapter~6). So $X^2-3$ is irreducible over $\mathbb F_M$ and
$K=\mathbb F_M(\sqrt3)$ is a field with $M^2$ elements (Tome~2, Chapter~9). Its Frobenius $x\mapsto x^M$ is a ring homomorphism fixing
$\mathbb F_M$, and $\sqrt3^M=3^{(M-1)/2}\sqrt3=-\sqrt3$ by Euler's criterion. Let $\tau=1+\sqrt3$; then $\tau^2=4+2\sqrt3=2\omega$ and
$\tau^{M+1}=\tau\cdot\tau^M=(1+\sqrt3)(1-\sqrt3)=-2$. Hence, with $\frac{M+1}2=2^{p-1}$,
\[
\omega^{2^{p-1}}=\Bigl(\frac{\tau^2}2\Bigr)^{(M+1)/2}=\frac{\tau^{M+1}}{2\cdot2^{(M-1)/2}}=\frac{-2}{2\bigl(\frac2M\bigr)}=-1.
\]
Then $s_{p-2}=\omega^{2^{p-2}}+\omega^{-2^{p-2}}=\omega^{-2^{p-2}}\bigl(\omega^{2^{p-1}}+1\bigr)=0$ in $K$. As $s_{p-2}$ is an integer, $M\mid s_{p-2}$.
\tools Tome~1, Chapter~6 (reciprocity, supplementary laws, Euler's criterion); Tome~2, Chapters~6 and~9 (quotient rings, finite
fields, Frobenius).
\leads Lucas (1876) proved $M_{127}$ prime by hand, a record that stood until computers; D.~H.~Lehmer (1930) gave the clean form.
Every record prime since 1952, except a few, is a Mersenne prime found by this test, now by the GIMPS project; the largest
known has more than $41$ million digits. The test is the analogue, in the group of norm-one elements of $\mathbb F_{M^2}$, whose
order $M+1$ is a power of $2$, of Pépin's test for Fermat numbers.
\end{solution}

\begin{solution}{pb:wolsten}
\idea Write $\binom{2p-1}{p-1}=\prod_{k=1}^{p-1}\frac{p+k}k=\prod\bigl(1+\frac pk\bigr)$. Modulo $p^3$ only the first three elementary symmetric
functions of the $\frac pk$ matter: $1+pH_{p-1}+p^2e_2$. So we need $H_{p-1}\equiv0\pmod{p^2}$ and $e_2\equiv0\pmod p$, and both come down
to the sum of the $\frac1{k^2}$, which is a sum of squares in disguise.

\writeup All fractions here have denominators prime to $p$. (a) As $k$ runs over $1,\dots,p-1$, so does the residue of $\frac1k$
(inversion is a bijection of $(\Z/p\Z)^*$). So $\sum\frac1{k^2}\equiv\sum_{k=1}^{p-1}k^2=\frac{(p-1)p(2p-1)}6\equiv0\pmod p$, since $6$ is
invertible modulo $p\ge5$.

(b) Pairing $k$ with $p-k$, $2H_{p-1}=\sum_{k=1}^{p-1}\bigl(\frac1k+\frac1{p-k}\bigr)=p\sum_k\frac1{k(p-k)}$. Modulo $p$,
$\frac1{k(p-k)}\equiv-\frac1{k^2}$, so the sum is $\equiv0\pmod p$ by (a). Hence $p^2\mid2H_{p-1}$, and $p^2\mid H_{p-1}$.

(c) $\binom{2p-1}{p-1}=\frac{(p+1)(p+2)\cdots(2p-1)}{(p-1)!}=\prod_{k=1}^{p-1}\Bigl(1+\frac pk\Bigr)=1+pe_1+p^2e_2+p^3(\cdots)$, with
$e_1=H_{p-1}$ and $e_2=\sum_{j<k}\frac1{jk}=\frac12\bigl(H_{p-1}^2-\sum\frac1{k^2}\bigr)$. By (b), $pe_1\equiv0\pmod{p^3}$; by (a) and (b),
$e_2\equiv0\pmod p$, so $p^2e_2\equiv0\pmod{p^3}$. Hence $\binom{2p-1}{p-1}\equiv1\pmod{p^3}$, and $\binom{2p}p=2\binom{2p-1}{p-1}\equiv2$.
\tools Tome~1, Chapter~3 (inverses modulo $p$); Tome~0, Chapters~1 and~3 (sums of squares, elementary symmetric functions).
\leads Wolstenholme (1862). Primes for which the congruence holds modulo $p^4$ are \emph{Wolstenholme primes}; only $16843$ and $2124679$ are
known. It is conjectured, and open, that the converse holds: if $\binom{2n-1}{n-1}\equiv1\pmod{n^3}$ then $n$ is prime.
\end{solution}

\begin{solution}{pb:jacobsthal}
\idea Split $ap$ objects into $(a-1)p$ and $p$, and use Vandermonde. The terms where the split is a multiple of $p$ give the
Pascal recursion for $\binom ab$; all other terms are products of two binomial coefficients each divisible by $p$.

\writeup Induction on $a$. For $a=b$ or $b=0$, both sides are $1$. Let $a>b\ge1$ (so $a\ge2$). By Vandermonde,
\[
\binom{ap}{bp}=\sum_i\binom{(a-1)p}i\binom p{bp-i},
\]
over $i$ with $0\le bp-i\le p$. If $p\mid i$, then $i=bp$ or $i=(b-1)p$, and these two terms are $\binom{(a-1)p}{bp}+\binom{(a-1)p}{(b-1)p}$. If
$p\nmid i$, then $0<bp-i<p$ and $p\mid\binom p{bp-i}$; also, with $N=(a-1)p$, $i\binom Ni=N\binom{N-1}{i-1}$ is divisible by $p$, and $p\nmid i$,
so $p\mid\binom Ni$. These terms are divisible by $p^2$. By induction,
$\binom{ap}{bp}\equiv\binom{a-1}b+\binom{a-1}{b-1}=\binom ab\pmod{p^2}$.
\tools Tome~0, Chapter~8 (Vandermonde, Pascal); Tome~1, Chapter~3 (binomial coefficients modulo $p$).
\leads Babbage (1819) proved $\binom{2p-1}{p-1}\equiv1\pmod{p^2}$, the case $a=2$, $b=1$; Wolstenholme then reached $p^3$. Jacobsthal showed that
$\binom{ap}{bp}\equiv\binom ab$ modulo $p^3$ for $p\ge5$, and even modulo $p^{3+v_p(ab(a-b))}$ (a direct computation verifies the $p^3$ version). Such
congruences are the $p$-adic continuity of the binomial coefficients, the beginning of $p$-adic analysis (Tome~12).
\end{solution}

\begin{solution}{pb:fermatquot}
\idea $a^{p-1}=1+pq_p(a)$: the Fermat quotient is the first digit of $a^{p-1}$ in base $p$ beyond the units digit, and
multiplying numbers of the form $1+px$ adds the $x$'s modulo $p$. So $q_p$ behaves like a logarithm. For (b), expand $2^p$ by the
binomial theorem; for (c), take the logarithm of Wilson's theorem.

\writeup All congruences are modulo $p$ unless stated. (a) $(ab)^{p-1}=(1+pq_p(a))(1+pq_p(b))\equiv1+p(q_p(a)+q_p(b))\pmod{p^2}$. Next,
$(a+pt)^{p-1}\equiv a^{p-1}+(p-1)pta^{p-2}\pmod{p^2}$, so $q_p(a+pt)\equiv q_p(a)-ta^{p-2}\equiv q_p(a)-\frac ta$, since $a^{p-1}\equiv1$.

(b) $2^p=\sum_k\binom pk=2+\sum_{k=1}^{p-1}\binom pk$, and $\frac1p\binom pk=\frac{(p-1)(p-2)\cdots(p-k+1)}{k!}\equiv\frac{(-1)^{k-1}(k-1)!}{k!}=\frac{(-1)^{k-1}}k$.
So $\frac{2^p-2}p\equiv\sum_{k=1}^{p-1}\frac{(-1)^{k-1}}k$, and $\frac{2^p-2}p=2q_p(2)$.

(c) Extend $q_p(N)=\frac{N^{p-1}-1}p$ to all $N$ prime to $p$; (a) still holds. Then
$\sum_{a=1}^{p-1}q_p(a)\equiv q_p\bigl((p-1)!\bigr)$. Write $(p-1)!=-1+pW_p$ (Wilson). As $p-1$ is even,
$(-1+pW_p)^{p-1}\equiv1-(p-1)pW_p\equiv1+pW_p\pmod{p^2}$, so $q_p((p-1)!)\equiv W_p$.

(d) Let $e$ and $E$ be the orders of $2$ modulo $p$ and $p^2$. Then $e\mid E$, and $2^{pe}=(1+cp)^p\equiv1\pmod{p^2}$ if $2^e=1+cp$, so
$E\mid pe$: $E\in\{e,pe\}$. If $2^{p-1}\equiv1\pmod{p^2}$, then $E\mid p-1$ is prime to $p$, so $E=e$. Conversely if $E=e$, then
$E\mid p-1$ and $2^{p-1}\equiv1\pmod{p^2}$.
\tools Tome~1, Chapters~3 and~5 (Wilson, Fermat, orders); Tome~0, Chapter~8 (binomial theorem).
\leads A prime with $q_p(2)\equiv0$ is a \emph{Wieferich prime}. Wieferich proved in 1909 that if $x^p+y^p=z^p$ has a solution with
$p\nmid xyz$, then $p$ is one. Only $1093$ (Meissner, 1913) and $3511$ (Beeger, 1922) are known; heuristically there should be about
$\log\log x$ of them below $x$, infinitely many but very few. Lerch's formula (c) connects them with Wilson primes ($W_p\equiv0$):
$5$, $13$ and $563$.
\end{solution}

\begin{solution}{pb:morley}
\idea $\binom{p-1}k=\prod_{j=1}^k\frac{p-j}j=(-1)^k\prod_j\bigl(1-\frac pj\bigr)\equiv(-1)^k(1-pH_k)\pmod{p^2}$. So the question is about the half
harmonic sum $H_{(p-1)/2}$ modulo $p$, and Problem~\ref{pb:fermatquot}(b) is exactly an alternating harmonic sum, which is
$H_{p-1}$ minus twice the even terms, that is minus $H_{(p-1)/2}$.

\writeup Let $k=\frac{p-1}2$. Expanding the product, $\binom{p-1}k\equiv(-1)^k(1-pH_k)\pmod{p^2}$. Next,
\[
\sum_{j=1}^{p-1}\frac{(-1)^{j-1}}j=H_{p-1}-2\sum_{i=1}^{k}\frac1{2i}=H_{p-1}-H_k\equiv-H_k\pmod p,
\]
because $H_{p-1}\equiv0\pmod p$ (pair $j$ and $p-j$ as in Problem~\ref{pb:wolsten}(b); this works for every odd $p$). By
Problem~\ref{pb:fermatquot}(b), $H_k\equiv-2q_p(2)\pmod p$. So $\binom{p-1}k\equiv(-1)^k(1+2pq_p(2))\pmod{p^2}$. On the other side,
$4^{p-1}=(1+pq_p(2))^2\equiv1+2pq_p(2)\pmod{p^2}$.
\tools Tome~1, Chapter~3; Problems~\ref{pb:wolsten} and~\ref{pb:fermatquot}.
\leads Morley (1895) proved the congruence modulo $p^3$ for $p\ge5$, which a direct computation verifies. Its sharpening is related to the
Bernoulli number $B_{p-3}$, and the numbers $\binom{p-1}{(p-1)/2}$ modulo $p^2$ appear in supercongruences for hypergeometric series
(``Ramanujan-type'' congruences), a lively research area.
\end{solution}

\begin{solution}{pb:glaisher}
\idea Sum the identity $k^{p-1}=1+pq_p(k)$ over $k=1,\dots,p-1$. Lerch's formula turns the right side into $p-1+pW_p$, that is
$(p-1)!+p$. The left side is a sum of powers, which Faulhaber's formula writes with Bernoulli numbers; modulo $p^2$ only the
term with $B_{p-1}$ survives, and von Staudt--Clausen controls the denominators.

\writeup By von Staudt--Clausen (Tome~4, Chapter~1), the denominator of $B_j$ ($j\ge2$ even) is the product of the primes $\ell$ with
$\ell-1\mid j$; so $B_j$ has no $p$ in its denominator for $j<p-1$, and $pB_{p-1}$ has none either. Faulhaber's formula with $n=p$
and exponent $p-1$ (Tome~4, Chapter~1) gives
\[
\sum_{k=0}^{p-1}k^{p-1}=\frac1p\sum_{j=0}^{p-1}\binom pjB_jp^{p-j}=\sum_{j=0}^{p-1}\binom pjB_jp^{p-1-j}.
\]
The term $j=p-1$ is $pB_{p-1}$. The term $j=0$ is $p^{p-1}\equiv0\pmod{p^2}$ ($p\ge3$). For $1\le j\le p-2$, $p\mid\binom pj$ and
$p\mid p^{p-1-j}$, and $B_j$ is $p$-integral, so these terms are $\equiv0\pmod{p^2}$. Hence $\sum_{k=1}^{p-1}k^{p-1}\equiv pB_{p-1}\pmod{p^2}$.

On the other hand, $\sum_{k=1}^{p-1}k^{p-1}=\sum_k(1+pq_p(k))=p-1+p\sum_kq_p(k)\equiv p-1+pW_p\pmod{p^2}$ by Lerch's formula
(Problem~\ref{pb:fermatquot}(c)). And $pW_p=(p-1)!+1$. So $(p-1)!+p\equiv pB_{p-1}\pmod{p^2}$.
\tools Tome~4, Chapter~1 (Faulhaber, von Staudt--Clausen); Tome~1, Chapter~3 (Wilson); Problem~\ref{pb:fermatquot}.
\leads Glaisher (1900). Reducing modulo $p$ gives back Wilson's theorem with $pB_{p-1}\equiv-1$, the $p$-part of von Staudt--Clausen. So
$p$ is a Wilson prime exactly when $pB_{p-1}\equiv p-1\pmod{p^2}$. The Bernoulli numbers modulo $p$ also decide whether $p$ is a
regular prime, Kummer's criterion for Fermat's last theorem (Tome~12).
\end{solution}

\begin{solution}{pb:gaussprimes}
\idea Tome~1 lists the Gaussian primes: $1+i$ (norm $2$), a conjugate pair of norm $p$ for each $p\equiv1\pmod4$, and the
primes $q\equiv3\pmod4$ themselves (norm $q^2$), each with four associates. Counting by norm reduces to counting primes in
progressions, and the asymptotics come from Tome~5.

\writeup (a) By Tome~1, Chapter~10, every Gaussian prime is associate to exactly one of: $1+i$; $\pi$ or $\bar\pi$, where
$p=\pi\bar\pi\equiv1\pmod4$ (and $\pi,\bar\pi$ are not associate); $q\equiv3\pmod4$ prime in $\Z$. Their norms are $2$, $p$, $q^2$, and
each class has $4$ elements (the units $\pm1,\pm i$). So $G(x)=4\bigl(1+2\pi(x;4,1)+\pi(\sqrt x;4,3)\bigr)$ for $x\ge2$.

(b) By the prime number theorem in progressions (Tome~5, Chapter~5), $\pi(x;4,1)\sim\frac x{2\log x}$, and
$\pi(\sqrt x;4,3)\le\sqrt x$. So $G(x)\sim\frac{4x}{\log x}$. The number of Gaussian integers of norm at most $x$ is the number of lattice
points in the disc of radius $\sqrt x$, $\sim\pi x$ (Tome~1, Chapter~10, or a direct count). By the same differencing as for
$\pi(x)$, the density of Gaussian primes near norm $N$ is $\frac4{\pi\log N}$: in the plane, a Gaussian integer of norm near $N$
is prime with probability about $\frac{4}{\pi}\cdot\frac1{\log N}$.

(c) The sum is $\frac12+2\sum_{p\le x,p\equiv1(4)}\frac1p+\sum_{q\equiv3(4),q^2\le x}\frac1{q^2}$. The last sum converges, and by
Dirichlet's theorem in Mertens form (Tome~5, Chapter~4, or by partial summation from the
$\frac{\log p}p$ version), $\sum_{p\le x,p\equiv1(4)}\frac1p=\frac12\log\log x+O(1)$. The total is $\log\log x+O(1)$.
\tools Tome~1, Chapter~10 (Gaussian primes); Tome~5, Chapters~4 and~5 (Dirichlet, prime number theorem in progressions).
\leads (c) is Mertens's theorem for the ring $\Z[i]$, and the analogue holds for every number field (Landau's prime ideal
theorem, Tome~12). The Gaussian primes are equidistributed in angle (Hecke, 1920), and Duke, Friedlander and Iwaniec showed that
there are Gaussian primes in very thin sectors. Whether one can walk to infinity on Gaussian primes with bounded steps (the
\emph{Gaussian moat problem}) is open.
\end{solution}

\begin{solution}{pb:mullin}
\idea Euclid's argument guarantees that each new term is a new prime. For (b), the number $P+1$ with $P=b_1\cdots b_n$ is prime to
$2$ and $3$ as soon as $3$ has appeared; if its largest prime factor were $5$, it would be a power of $5$, and $5^c-1$ is
divisible by $4$, whereas $P$ contains the prime $2$ only once.

\writeup (a) $a_{n+1}$ divides $a_1\cdots a_n+1$, which is prime to each $a_i$, so $a_{n+1}\ne a_i$. The products give:
$2+1=3$; $6+1=7$; $42+1=43$; $1806+1=13\cdot139$, so $a_5=13$; $23478+1=23479=53\cdot443$, so $a_6=53$;
$23478\cdot53+1=1244335=5\cdot248867$, so $a_7=5$.

(b) The terms are distinct for the same reason, so for $n\ge1$, $P_n=b_1\cdots b_n$ contains the factor $2$ exactly once. $b_2=3$. For
$n\ge2$, $6\mid P_n$, so $P_n+1$ is prime to $2$ and $3$. If its largest prime factor were $5$, then $P_n+1=5^c$ with $c\ge1$, and
$P_n=5^c-1\equiv0\pmod4$, contradicting $v_2(P_n)=1$. And $b_2=3\ne5$. So $5$ never appears.
\tools Tome~1, Chapters~1 and~2 (Euclid's argument, factorisation).
\leads Mullin asked in 1963 whether every prime appears in the first sequence; Shanks conjectured yes, and it is open, the
sequence being known for only about fifty terms because each one needs a hard factorisation. For the second sequence, Cox and
van der Poorten (1968) showed that $5,11,13,17,19,23,29,31,37,41,47,53$ never appear, by arguments like (b) with more
congruences, and Booker (2012) proved that it misses infinitely many primes.
\end{solution}

\begin{solution}{pb:nsq1}
\idea This is Chebyshev's idea: compare the size of a product with the largest possible contribution of small primes. The
product $\prod(n^2+1)$ has logarithm about $2x\log x$. A prime $p\le x$ can divide $n^2+1$ only if $p=2$ or $p\equiv1\pmod4$, and
then for at most two residues $n$ modulo each $p^k$. So the primes $p\le x$ contribute about $2x\sum_{p\equiv1(4)}\frac{\log p}p\approx x\log x$:
half of what is needed. The missing half must come from primes larger than $x$.

\writeup Let $Q=\prod_{n\le x}(n^2+1)$. Then $\log Q\ge\sum_{n\le x}2\log n=2\log\lfloor x\rfloor!=2x\log x+O(x)$ (Tome~0, Chapter~5 or
Stirling). Suppose that every prime factor of $Q$ is at most $x$. If $p\mid n^2+1$, then $p=2$ or $p\equiv1\pmod4$ (Tome~1, Chapter~6).
For $p=2$, $v_2(n^2+1)\le1$ (squares are $0$ or $1$ modulo $4$), so $v_2(Q)\le x$. For odd $p$ and $k\ge1$, $n^2\equiv-1\pmod{p^k}$
has at most $2$ solutions modulo $p^k$ (the group $(\Z/p^k\Z)^*$ is cyclic, and $y^2=c$ has at most two solutions in a cyclic
group), so at most $2\bigl(\frac x{p^k}+1\bigr)$ solutions with $n\le x$, and none if $p^k>x^2+1$. So
\[
v_p(Q)=\sum_{k\ge1}\#\{n\le x:p^k\mid n^2+1\}\le\frac{2x}{p-1}+\frac{2\log(x^2+1)}{\log p}.
\]
Summing $v_p(Q)\log p$ over the primes $p\le x$ with $p\equiv1\pmod4$, and adding $x\log2$ for $p=2$:
\[
\log Q\le x\log2+2x\sum_{\substack{p\le x\\p\equiv1(4)}}\frac{\log p}{p-1}+2\pi(x)\log(x^2+1).
\]
By Dirichlet's theorem (Tome~5, Chapter~4),
\[
\sum_{p\le x,\ p\equiv1\ (4)}\frac{\log p}p=\frac12\log x+O(1),
\]
and
$\frac{\log p}{p-1}-\frac{\log p}p=\frac{\log p}{p(p-1)}$ has a convergent sum. By Chebyshev (Tome~1, Chapter~2), $\pi(x)\log x=O(x)$. So
$\log Q\le x\log x+O(x)$, which contradicts $\log Q\ge2x\log x+O(x)$ for large $x$. Hence $P(x)>x$ for all large $x$.
\tools Tome~1, Chapters~2, 5 and~6 (Chebyshev, cyclic groups, $-1$ as a square); Tome~5, Chapter~4 (Dirichlet's theorem).
\leads Chebyshev stated, and Markov proved (1895), that $P(x)/x\to\infty$. Deshouillers and Iwaniec (1982) showed that $n^2+1$ has a
prime factor larger than $n^{1.2}$ infinitely often. Whether $n^2+1$ is itself prime infinitely often is one of Landau's four
problems (1912), all still open.
\end{solution}

\begin{solution}{pb:nair}
\idea A polynomial integral $\int_0^1x^{m-1}(1-x)^{n-m}\dd x$, expanded, is a sum of fractions $\frac{\pm c}{m+k}$ with denominators at most
$n$, so $d_n$ times it is an integer. Computed with the Beta function, it is $\frac1{m\binom nm}$. So $m\binom nm\mid d_n$, and the central
binomial coefficient gives the exponential lower bound.

\writeup (a) Expanding $(1-x)^{n-m}$,
\[
I=\int_0^1x^{m-1}(1-x)^{n-m}\dd x=\sum_{k=0}^{n-m}\binom{n-m}k\frac{(-1)^k}{m+k},
\]
and $m+k\le n$, so $d_nI\in\Z$. By the Beta integral (Tome~4, Chapter~5), $I=\frac{(m-1)!(n-m)!}{n!}=\frac1{m\binom nm}$. So $m\binom nm\mid d_n$.

(b) With $n=2k$, $m=k$: $k\binom{2k}k\mid d_{2k}\mid d_{2k+1}$. With $n=2k+1$, $m=k+1$: $(k+1)\binom{2k+1}{k+1}=(2k+1)\binom{2k}k$ divides
$d_{2k+1}$. As $\gcd(k,2k+1)=1$, $k(2k+1)\binom{2k}k\mid d_{2k+1}$. The largest of the $2k+1$ terms $\binom{2k}j$ of $(1+1)^{2k}$ is at
least their mean: $\binom{2k}k\ge\frac{4^k}{2k+1}$. So $d_{2k+1}\ge k\,4^k$. For $k\ge2$, $k4^k\ge2^{2k+1}$: $d_n\ge2^n$ for odd $n\ge5$. For
even $n=2k+2$, $d_n\ge d_{2k+1}\ge k4^k\ge2^{2k+2}$ as soon as $k\ge4$, that is $n\ge10$; and $d_8=840\ge256$. Together with
$d_7=420\ge128$ and $d_9=2520\ge512$, $d_n\ge2^n$ for all $n\ge7$.

Finally $d_n=\prod_{p\le n}p^{\lfloor\log n/\log p\rfloor}\le n^{\pi(n)}$, since each $p^{\lfloor\log n/\log p\rfloor}\le n$. So
$\pi(n)\log n\ge\log d_n\ge n\log2$.
\tools Tome~4, Chapter~5 (Beta function); Tome~1, Chapters~1 and~2 (lcm, Chebyshev bounds); Tome~0, Chapter~8 (binomial coefficients).
\leads Nair (1982). Note that $\log d_n=\psi(n)$, Chebyshev's function, so the prime number theorem is equivalent to
$d_n=e^{n(1+o(1))}$. The same integrals, with more variables, are the start of Beukers's proofs of the irrationality of $\zeta(2)$ and
$\zeta(3)$ (Chapter~9), where the bound $d_n\le3^n$ for large $n$ is exactly what is needed.
\end{solution}

\begin{solution}{pb:artin}
\idea The group $(\Z/p\Z)^*$ has order $2q$, so the order of $2$ is $1$, $2$, $q$ or $2q$. Orders $1$ and $2$ are excluded by size. Order
$q$ means $2^q=2^{(p-1)/2}\equiv1$, which by Euler's criterion means that $2$ is a square; and whether $2$ is a square depends on $p$
modulo $8$, which is fixed by $q$ modulo $4$.

\writeup The order $e$ of $2$ divides $p-1=2q$. $e=1$ or $2$ would need $p\mid1$ or $p\mid3$, excluded. If $q\equiv1\pmod4$, $p=2q+1\equiv3\pmod8$,
so $\bigl(\frac2p\bigr)=-1$ and $2^q\equiv-1$ by Euler's criterion: $e\ne q$, so $e=2q$ and $2$ is a primitive root. If $q\equiv3\pmod4$, then
$p\equiv7\pmod8$ and $\bigl(\frac2p\bigr)=1$, $\bigl(\frac{-1}p\bigr)=-1$, so $\bigl(\frac{-2}p\bigr)=-1$; the same argument applies to $-2$, whose order is
not $1$ or $2$ since $p>3$.
\tools Tome~1, Chapters~5 and~6.
\leads Artin conjectured (1927) that $2$ is a primitive root modulo infinitely many primes, in fact for a proportion $\prod_p\bigl(1-\frac1{p(p-1)}\bigr)
=0.3739\ldots$ of them. Hooley (1967) proved it under the generalised Riemann hypothesis. Unconditionally, Heath-Brown (1986) showed that at most two
primes fail to be primitive roots for infinitely many $p$: so one of $2$, $3$, $5$ works, but nobody knows which. This problem would
give a proof if there were infinitely many such ``safe primes'', which is also open.
\end{solution}

\begin{solution}{pb:findprime}
\idea A draw succeeds with probability $\frac{\pi(2N)-\pi(N)}{N+1}$, and the number of draws until the first success is geometric, with
mean the inverse of that. The prime number theorem gives the probability; sieving out small factors raises it by the factor by
which they thin out the candidates.

\writeup (a) By the prime number theorem (Tome~5, Chapter~5), $\pi(2N)-\pi(N)\sim\frac{2N}{\log2N}-\frac N{\log N}\sim\frac N{\log N}$. The
number of independent uniform draws until a prime is geometric with parameter $p_N\sim\frac1{\log N}$, of mean $\frac1{p_N}\sim\log N$
(Tome~3, Chapter~2).

(b) The integers of $[N,2N]$ prime to $30$ number $\frac8{30}N+O(1)$, and all primes $>5$ are among them. So the success probability
becomes $\sim\frac{N/\log N}{(4/15)N}$, and the mean number of draws $\sim\frac4{15}\log N$. With all primes $\le y$, the factor is
$\prod_{p\le y}(1-\frac1p)$, which by Mertens's theorem (Tome~5, Chapter~1) is $\sim\frac{e^{-\gamma}}{\log y}$: sieving to $y$ divides the
work by about $1.78\log y$.

(c) $\log2^{1023}=1023\log2\approx709$, so about $709$ draws; with the wheel $30$, about $189$. Each draw that survives costs one
Miller--Rabin test, the others a few trial divisions.
\tools Tome~5, Chapters~1 and~5 (Mertens, prime number theorem); Tome~3, Chapter~2 (geometric law).
\leads Real implementations sieve an interval rather than draw independently, which biases the choice: primes after a long gap are
chosen more often. Brandt and Damgård (1993) and later Fouque and Tibouchi analysed this bias and its (small) effect on security.
\end{solution}

\begin{solution}{pb:leastqnr}
\idea If $n=n_p$ were composite, its factors, being smaller, would be residues, and so would $n$. For the size, use multiples:
the first multiple of $n$ beyond $p$ is $mn$, and $mn-p<n$ is a residue, so $m$ must be a non-residue, hence $m\ge n$. That forces
$n^2\lesssim p$. In the random model, $n_p=p_k$ when the first $k-1$ prime symbols are $+1$ and the $k$-th is $-1$.

\writeup (a) Every $1\le a<n_p$ is a residue. If $n_p=ab$ with $1<a,b<n_p$, then $\bigl(\frac{n_p}p\bigr)=\bigl(\frac ap\bigr)\bigl(\frac bp\bigr)=1$,
absurd; so $n_p$ is prime. Let $n=n_p$ and $m=\lfloor\frac pn\rfloor+1$, so $0<mn-p<n$ (here $n\nmid p$, because $1<n<p$: the
number $1$ is a residue, and non-residues exist in $\{1,\dots,p-1\}$). Then $mn-p$ is a residue, so
$\bigl(\frac mp\bigr)\bigl(\frac np\bigr)=\bigl(\frac{mn}p\bigr)=\bigl(\frac{mn-p}p\bigr)=1$, $\bigl(\frac mp\bigr)=-1$, and $m\ge n$. So $n\le\frac pn+1$,
$n^2-n\le p$, and $n\le\frac{1+\sqrt{1+4p}}2<\sqrt p+1$.

(b) In the model, $P(n_p=p_k)=2^{-(k-1)}\cdot\frac12=2^{-k}$, so $E[n_p]=\sum_kp_k2^{-k}=\frac22+\frac34+\frac58+\frac7{16}+\dots=3.6746\ldots$.
\tools Tome~1, Chapter~6 (Legendre symbol); Tome~3, Chapter~2 (expectation of a discrete variable).
\leads Erdős (1961) proved that the average of $n_p$ over primes $p\le x$ tends to $\sum p_k2^{-k}$, so the model is right on average; the
one finds numerically $3.61$ for $p<2\cdot10^5$. The worst case is another story: Burgess (1957) showed $n_p\ll p^{1/(4\sqrt e)+\varepsilon}$, still the
best known exponent, while the generalised Riemann hypothesis gives $n_p<2\log^2p$ (Ankeny, Bach), which is what makes Miller's
deterministic test work.
\end{solution}

\section*{Map of the chapter}

\begin{center}
\small
\adjustbox{max width=\linewidth}{\begin{tabular}{lll}
Problem & Topic & Tomes used\\\hline
\ref{pb:fermatnum} & Fermat numbers & 1\\
\ref{pb:mersennediv} & divisors of Mersenne numbers & 1\\
\ref{pb:pseudo} & pseudoprimes & 1\\
\ref{pb:carmichael} & Carmichael numbers & 1, 2\\
\ref{pb:euler} & Solovay--Strassen & 1, 2\\
\ref{pb:aks} & the AKS identity & 1\\
\ref{pb:pratt} & Pratt certificates & 1\\
\ref{pb:pocklington} & Pocklington and Proth & 1\\
\ref{pb:rabin} & Miller--Rabin, one in four & 1, 2\\
\ref{pb:lucaslehmer} & Lucas--Lehmer & 1, 2\\
\ref{pb:wolsten} & Wolstenholme & 1, 0\\
\ref{pb:jacobsthal} & $\binom{ap}{bp}$ modulo $p^2$ & 0, 1\\
\ref{pb:fermatquot} & Fermat quotients, Lerch & 1, 0\\
\ref{pb:morley} & Morley's congruence & 1\\
\ref{pb:glaisher} & Glaisher, Bernoulli & 4, 1\\
\ref{pb:gaussprimes} & Gaussian primes & 1, 5\\
\ref{pb:mullin} & Euclid--Mullin & 1\\
\ref{pb:nsq1} & prime factors of $n^2+1$ & 1, 5, 0\\
\ref{pb:nair} & lcm and Chebyshev & 4, 1\\
\ref{pb:artin} & primitive roots for safe primes & 1\\
\ref{pb:findprime} & drawing random primes & 5, 3\\
\ref{pb:leastqnr} & least non-residue & 1, 3\\
\end{tabular}}
\end{center}

\section*{Sources and further reading}

\begin{itemize}
\item F.~N.~Cole, On the factoring of large numbers, \emph{Bulletin of the American Mathematical Society} 10 (1903), 134--137; E.~T.~Bell,
\emph{Mathematics: Queen and Servant of Science}, McGraw-Hill, 1951, p.~228 (the story of the silent lecture and of
the ``three years of Sundays''; the \emph{Bulletin} note is the only contemporary record of the talk).
\item R.~Crandall and C.~Pomerance, \emph{Prime Numbers: A Computational Perspective}, 2nd ed., Springer, 2005. Pseudoprimes,
Carmichael numbers, Lucas--Lehmer, Pocklington, Proth, Miller--Rabin, AKS and much more.
\item W.~R.~Alford, A.~Granville and C.~Pomerance, There are infinitely many Carmichael numbers, \emph{Annals of Mathematics} 139 (1994),
703--722.
\item M.~O.~Rabin, Probabilistic algorithm for testing primality, \emph{Journal of Number Theory} 12 (1980), 128--138; L.~Monier,
Evaluation and comparison of two efficient probabilistic primality testing algorithms, \emph{Theoretical Computer Science} 12 (1980),
97--108.
\item V.~R.~Pratt, Every prime has a succinct certificate, \emph{SIAM Journal on Computing} 4 (1975), 214--220.
\item M.~Agrawal, N.~Kayal and N.~Saxena, PRIMES is in P, \emph{Annals of Mathematics} 160 (2004), 781--793.
\item A.~Granville, Binomial coefficients modulo prime powers, in \emph{Organic Mathematics}, CMS Conference Proceedings 20 (1997), 253--275.
Wolstenholme, Babbage, Jacobsthal and Morley in one place.
\item R.~J.~McIntosh, On the converse of Wolstenholme's theorem, \emph{Acta Arithmetica} 71 (1995), 381--389.
\item M.~Lerch, Zur Theorie des Fermatschen Quotienten, \emph{Mathematische Annalen} 60 (1905), 471--490; R.~Crandall, K.~Dilcher and
C.~Pomerance, A search for Wieferich and Wilson primes, \emph{Mathematics of Computation} 66 (1997), 433--449.
\item A.~A.~Mullin, Recursive function theory, \emph{Bulletin of the American Mathematical Society} 69 (1963), 737; C.~D.~Cox and
A.~J.~van der Poorten, On a sequence of prime numbers, \emph{Journal of the Australian Mathematical Society} 8 (1968), 571--574;
A.~R.~Booker, On Mullin's second sequence of primes, \emph{Integers} 12 (2012), 1167--1177.
\item J.-M.~Deshouillers and H.~Iwaniec, On the greatest prime factor of $n^2+1$, \emph{Annales de l'Institut Fourier} 32 (1982), 1--11.
\item M.~Nair, On Chebyshev-type inequalities for primes, \emph{American Mathematical Monthly} 89 (1982), 126--129.
\item C.~Hooley, On Artin's conjecture, \emph{Journal für die reine und angewandte Mathematik} 225 (1967), 209--220; D.~R.~Heath-Brown,
Artin's conjecture for primitive roots, \emph{Quarterly Journal of Mathematics} 37 (1986), 27--38.
\item P.~Erdős, Remarks on number theory I, \emph{Matematikai Lapok} 12 (1961), 10--17 (the average least non-residue); D.~A.~Burgess, The
distribution of quadratic residues and non-residues, \emph{Mathematika} 4 (1957), 106--112.
\end{itemize}
